\chapter*{Introduzione}
\addcontentsline{toc}{chapter}{Introduzione}

Recentemente le simulazioni numeriche sono diventate un requisito fondamentale 
per la progettazione di diverse componenti ingegneristiche. 
Infatti in fase di progettazione di sistemi complessi, in particolar modo in 
ambito nucleare, è richiesto lo sviluppo di codici di calcolo differenti in 
grado di garantire una simulazione numerica del progetto, che sia attendibile e 
inerente ai processi fisici reali.
Per questa ragione molti problemi multifisici e multiscala sono studiati 
attraverso l'accoppiamento di codici differenti appartenenti a discipline diverse che 
risolvono i singoli problemi fisici.

In questo progetto di tesi è stato implementato un modello di reattore nucleare 
LFR (Lead Fast Reactor) all'interno di una piattaforma numerica, la quale 
integra differenti strumenti computazionali allo scopo di investigare il 
comportamento multifisico ottenuto dall'accoppiamento di codici di neutronica e 
di termoidraulica.
Il modello di reattore nucleare formalizzato è basato sul progetto ALFRED 
(Advanced Lead Fast REactor Demonstrator), il quale ha il ruolo di dimostrare la 
fattibilità della tecnologia dei reattori veloci raffreddati al piombo proposti 
dal GIF (Generation IV International Forum) come sistemi energetici nucleari 
appartenenti alla IV$^{a}$ generazione.
I reattori veloci raffreddati al piombo dispongono di caratteristiche innovative di 
sostenibilità e di affidabilità offerte dalle proprietà del piombo come 
refrigerante; quest'ultimo è in grado di mantenere uno spettro neutronico veloce 
che massimizza la fertilizzazione del combustibile, prospettando la possibilità di 
ottenere un ciclo chiuso di combustibile.

La possibilità di studiare sistemi nucleari attraverso simulazioni numeriche 
presenta molteplici vantaggi durante la progettazione. Inoltre  l'integrazione
di codici di neutronica e di termoidraulica all'interno di una 
piattaforma numerica permette di scambiare dati tra essi e di assicurare 
un approccio multifisico completo.
In particolare, all'interno di questa tesi è presentato l'accoppiamento del 
codice \textit{Donjon (Full-core simulation)}, che si occupa della simulazione 
neutronica dell'intero reattore e del codice \textit{FEMuS}, che risolve la
termoidraulica, all'interno della piattaforma numerica.
La piattaforma numerica, basata su \textit{Salom\'e}, permette la gestione e lo 
scambio delle strutture dati offerte da librerie e moduli. In particolare, 
la distribuzione di potenza termica ottenuta dalla risoluzione del flusso 
neutronico all'interno del codice di reattore viene passata al codice di 
termoidraulica al fine di risolvere il campo di temperatura sul dominio del 
reattore.
La distribuzione di temperatura è trasferita, a sua volta, al codice di 
neutronica al fine di interpolare sul database reattoristico e di ottenere le 
sezioni d'urto omogeneizzate e condensate a 33 gruppi energetici in funzione 
della temperatura del combustibile, della temperatura del refrigerante e dei tempi di 
bruciamento isotopico del combustibile fissile.
Il database reattoristico è stato realizzato attraverso il supporto del codice 
di cella \textit{Dragon (Lattice code)}, il quale risolve il problema del 
trasporto dei neutroni ottenendo le sezioni d'urto efficaci relative alle 
principali componenti del reattore come le barre di combustibile, le barre di 
controllo, elementi di protezione neutronica. I dati delle librerie nucleari 
associati alle sezioni d'urto microscopiche di ogni elemento sono stati estratti 
dalle \textit{Jeff3.1.2shem315}.

Attraverso la versatilità della piattaforma numerica è stato possibile realizzare 
un ciclo iterativo tra i codici di neutronica e di termoidraulica consentendo il calcolo del
% l'analisi multifisica del modello di ALFRED e la stampa della distribuzione del 
flusso neutronico, della potenza termica e della temperatura.
L'analisi neutronica del reattore è stata investigata attraverso la valutazione 
del fattore moltiplicativo effettivo $k_{eff}$ e del coefficiente di reattività 
$\Delta \rho$ in funzione delle variazioni di temperatura, che incidono sulle 
proprietà nucleari dei materiali e del bruciamento isotopico, considerando il
tempo di inizio ciclo (BOC) di 2 anni e il tempo di fine ciclo di 3 anni.
Il campo di temperatura è stato ottenuto attraverso la risoluzione 
dell'equazione di conservazione dell'energia ipotizzando la velocità del 
refrigerante costante e il coefficiente di conducibilità termica mediato sul 
piombo.

% Sono stati ottenuti valori del fattore moltiplicativo effettivo compatibili con 
% altri studi inerenti al progetto ALFRED; persino la temperatura di uscita del 
% refrigerante ottenuta calcolando la potenza termica neutronica è risultata accettabile 
% con i dati progettuali previsti.
% I risultati ottenuti soddisfano i parametri previsti per il progetto ALFRED; tuttavia sarebbe necessaria la raffinazione dei dati utilizzati durante la 
% modellazione del reattore in riferimento ai cambiamenti apportati dagli enti di 
% ricerca che operano nella costruzione del progetto dimostrativo ALFRED. 
% Inoltre è auspicabile l'introduzione di modelli più accurati e complessi nell'analisi termoidraulica del reattore, 
% come modelli per descrivere la turbolenza relativa al piombo liquido.

% In conclusione la possibilità di integrare l'accoppiamento 
% neutronico-termoidraulica in un unica piattaforma numerica può essere un ottimo 
% strumento per la progettazione e la valutazione della fattibilità della 
% tecnologia LFR, aprendo a prospettive future per indagare il comportamento 
% multifisico di sistemi energetici complessi. 




Nel Capitolo \ref{cap1} sono trattate le tematiche relative ai reattori veloci 
di IV$^{a}$ generazione in particolare basati sulla tecnologia di raffreddamento al 
piombo LFR.
I reattori veloci hanno caratteristiche uniche di affidabilità e di 
sostenibilità, in grado di ottenere un ciclo chiuso di combustibile con 
attenzione ai problemi ambientali. Inoltre è approfondito il progetto europeo 
ALFRED, il cui modello è stato utilizzato per valutare l'accoppiamento 
multifisico della piattaforma numerica.
Nel Capitolo \ref{cap2} sono affrontate le principali equazioni della neutronica 
e la loro discretizzazione al fine di essere risolte numericamente attraverso 
tecniche e metodi implementati nei codici di calcolo neutronico, sia di cella 
che di intero reattore.
Nel Capitolo \ref{cap3} sono derivate le equazioni della termodinamica di 
conservazione della massa, della quantità di moto e dell'energia dall'equazione 
di conservazione per una grandezza estensiva. Sono accennate le principali 
approssimazioni di termoidraulica per il reattore nucleare.
Nel Capitolo \ref{cap4} sono presentati i codici di neutronica e la loro 
implementazione nella piattaforma numerica.
Viene descritto il procedimento per la realizzazione del database reattoristico, 
ottenuto dal codice di cella, e la successiva interpolazione sulle sezioni 
d'urto efficaci per la risoluzione del problema neutronico associato al modello 
di reattore ALFRED. 
Infine, nel Capitolo \ref{cap5} sono presentati e valutati i risultati ottenuti 
dall'accoppiamento neutronico-termoidraulico in relazione al reattore considerato. 


